% Rules at any angle: a square root, a fraction, a QED box and a frame, upright and in \rotatebox.
\documentclass{article}
\usepackage{amsmath,amsthm,graphicx}
\begin{document}
Rules are the bars of fractions such as $\frac{a+b}{2}$, the vinculum of a
square root like $\sqrt{x^2+y^2}$, and the four sides of the box that ends
a proof. Turned by \rotatebox{45}{$\sqrt{x+1}$} forty-five degrees, each
is a rule no longer upright, and so is a frame turned by thirty,
\rotatebox{30}{\fbox{framed}}, or a whole fraction turned the other way,
\rotatebox{-45}{$\frac{1}{n}$}, in the middle of a sentence that goes on
long enough to break across lines differently at every width.
\begin{proof}
The square root of two is not a fraction $p/q$, for then $p^2 = 2q^2$ and
both would be even; so its expansion $\sqrt{2} = 1.41421\ldots$ never ends,
and the proof ends with a box turned on its corner.
\renewcommand{\qedsymbol}{\rotatebox{45}{\openbox}}
\end{proof}
\begin{proof}
And one upright, for comparison, after a sentence of its own.
\end{proof}
\end{document}
